\documentclass[11pt,article,oneside]{memoir}
\usepackage[utf8]{inputenc}
\usepackage[T1]{fontenc}
\usepackage{url}
\usepackage[hidelinks,colorlinks=true,linkcolor=blue]{hyperref}
\usepackage{amsmath,amssymb}
\usepackage{graphicx}
\usepackage{xcolor}
\usepackage{tabularx}
\usepackage{listings}
\usepackage{geometry}
\usepackage{titlesec}
\usepackage{libertine}

\geometry{margin=1in}

\titleformat{\section}{\Large\bfseries\color{darkgray}}{\thesection}{1em}{}
\titleformat{\subsection}{\large\bfseries\color{blue!80!black}}{\thesubsection}{1em}{}

\lstset{
  basicstyle=\small\ttfamily,
  breaklines=true,
  frame=lines,
  backgroundcolor=\color{gray!10}
}

\begin{document}

\begin{center}
    \Huge \textbf{Plain Language Guide Professional} \\
    \vspace{0.5em}
    \large \textit{A CEREBRUM Research Publication (v1.2.4)} \\
    \vspace{1em}
    \large \textbf{Bryan Alexander Buchorn} \\
    Independent Researcher \\
    \vspace{2em}
\end{center}

\begin{abstract}
\textbf{Overview:} This guide provides a conceptual entry point into the CEREBRUM framework. It is designed for practitioners and researchers seeking to understand the core reasoning signals of Framework Implementation Guide.
\end{abstract}

\textless div class="cover"\textgreater 
    \textless img src="file:///C:/Users/bryan/.gemini/antigravity/brain/77bb37a0-e733-41be-824d-b07e7cce5a6f/cerebrum\_hero\_v77bb37a0\_e730\_41be\_824d\_b07e7cce5a6f\_png\_1774651388694.png" alt="CEREBRUM Visionary Hero"\textgreater 
    \textless h1 class="title"\textgreater CEREBRUM EXPLAINED\textless /h1\textgreater 
    \textless p class="subtitle"\textgreater A Plain-Language Guide to the Research\textless /p\textgreater 
    \textless div class="meta"\textgreater 
        \textless strong\textgreater Version 3.0 (v1.2.4 Hardened Edition)\textless /strong\textgreater \textless br\textgreater 
        March 2026 — Independent Researcher
    \textless /div\textgreater 
\textless /div\textgreater 

---

\subsubsection{\textbf{Who this is for}}
You're comfortable with computers and technology, but you don't have a background in machine learning or knowledge graphs. This guide explains CEREBRUM's individual pieces, name each piece in plain English, and walk through how the whole system works.

---

\subsubsection{\textbf{How to Read This Guide}}
Each chapter builds on the last. If you're already comfortable with a concept, feel free to skip ahead. Every technical term is explained when it first appears. Formulas appear in clearly marked "Math, Unpacked" sections — you can read around them without losing the narrative, or dive in for a full under\-standing.

\begin{center}
{\footnotesize
\begin{tabularx}{\columnwidth}{|>{\raggedright\arraybackslash}X|>{\raggedright\arraybackslash}X|>{\raggedright\arraybackslash}X|}
\hline
Chapter & Topic & You'll Understand \\ \hline
1 & What is a Knowledge Graph? & The database behind CEREBRUM \\ \hline
2 & What is an LLM? & What ChatGPT and its peers actually are \\ \hline
3 & The problem with current AI reasoning & Why we need something new \\ \hline
4 & What is attention in AI? & The mechanism CEREBRUM borrows \\ \hline
5 & What is community detection? & The other key concept \\ \hline
6 & What is CEREBRUM? & The full picture — including THALAMUS and CORTEX \\ \hline
7 & Why does it matter? & Real-world stakes \\ \hline
8 & What are we trying to prove? & The research questions \\ \hline
9 & Does it work? & Validation results (v1.2.4) \\ \hline
10 & The Studio: Seeing it live & The interactive interface \\ \hline
11 & Real-time streaming data & Sensors, signals, and live feeds \\ \hline
12 & What kinds of data can it use? & Beyond sprea\-dsheets \\ \hline
13 & Bayesian Beam Search & Reasoning under uncertainty \\ \hline
14 & The REM Cycle & Self-maintenance and synthesis \\ \hline
15 & Insight Learning & When the graph has a Eureka moment \\ \hline
16 & The Formal Journey & ArXiv, Margins, and Academic Rigor \\ \hline
\end{tabularx}
}
\end{center}

---

\subsection{CHAPTER 1 — What Is a Knowledge Graph?}
Most databases are like sprea\-dsheets: rows and columns. A \textbf{Knowledge Graph} is different. It's a collection of entities (nodes) and the relat\-ionships (edges) between them. Instead of a row for "Tom Hanks," a Knowledge Graph has a circle for \textbf{Tom Hanks}. That circle is connected by a line labeled \textbf{ACTED\_IN} to another circle labeled \textbf{Forrest Gump}. This format is how humans think — we think in connections, not table rows.

---

\subsection{CHAPTER 3 — The Problem With AI Reasoning}
Because LLMs are statistical, they "hallucinate." If they don't know an answer, they'll often make one up that \textit{sounds} perfectly plausible. \textbf{CEREBRUM is different.} It moves the reasoning out of the proba\-bilistic "Black Box" and into the graph itself. It only follows lines that actually exist in your data.

---

\subsection{CHAPTER 9 — Does It Work? (v1.2.4 Results)}
As of late March 2026, CEREBRUM has been validated on several large-scale datasets and is operating at \textbf{v1.2.4 "Camera-Ready"} status.

\textbf{What the tests showed:}
\begin{enumerate}
\item \textbf{Unrivaled Stability}: The framework now passes \textbf{1,042 automated tests}, including 130+ advanced "edge-case" scenarios.
\item \textbf{Medical accuracy}: On the Hetionet biomedical graph, CEREBRUM was over \textbf{183\% more accurate} than simple search at finding connections between diseases and genes.
\item \textbf{Reasoning Recall}: On MetaQA 3-hop reasoning, recall improved by \textbf{+350\%} relative to non-attention baselines.
\item \textbf{Zero training required}: Most AI systems require weeks of expensive training. CEREBRUM achieved these results with no training at all — it reasons purely from the structure of the graph it's given.

\end{enumerate}
---

\subsection{CHAPTER 16 — The Formal Journey: ArXiv and Beyond}
The transition from a laboratory project to a formal \textbf{ArXiv Manuscript} required a process called \textbf{Hardening}. This included:
\begin{itemize}
\item \textbf{Margin Remediation}: Ensuring every equation and table fits perfectly into a profe\-ssional two-column academic layout.
\item \textbf{Structural Hole Repair}: Identifying 12 subtle failure modes (like "Zombie Bridges" or "Causal Floods") and building mathe\-matical guards against them.
\item \textbf{Branding Audit}: Ensuring absolute privacy and profe\-ssional attribution across thousands of lines of code.

\end{itemize}
---

\subsubsection{\textbf{Acknow\-ledgments \& Credits}}
CEREBRUM stands on the shoulders of decades of found\-ational research in Graph Theory, Neuros\-cience, and Machine Learning.

\textbf{Questions? Contact the author:}
\textbf{Bryan Alexander Buchorn} — bryan.alexander@buchorn.com

---
\textit{Generated: March 2026 — Project CEREBRUM v1.2.4 — Final Rebranding}


\vspace{3em}
\begin{center}
    \small \textit{Project CEREBRUM — March 2026 — Hardened Enterprise Security (v1.2.4)}
\end{center}

\end{document}
